\section{A native, source-checked compact knot exterior}
\label{sec:diagram-exterior}

The normal-surface routines of Sections~\ref{sec:weighted-components} and
\ref{sec:compact-coordinates} accept a supplied compact manifold. Their
manifold checks alone do not establish that it is the exterior of the input
knot. We now supply that missing construction for one canonical, deliberately
unsimplified triangulation. The native producer starts from the validated PD,
uses four crossing cells per crossing, and replaces each by five finite
tetrahedra. A separate checker recovers the same gluing specification from
the PD and integer barycentric coordinates. Neither routine imports Regina.
Normal-vector discovery and verified simplification remain further tasks.

\subsection{Crossing cells and the source theorem}

Let $n$ be the number of crossings and put $N=\max(1,n)$. The zero-crossing
circle is represented by the positive Reidemeister-I curl with PD
$(0,1,1,0)$. Otherwise use the original normalized PD. Its four slots are
counterclockwise, slots $0,2$ belong to the under-strand, and slots $1,3$
belong to the over-strand. The existing validator checks the spherical
rotation system and the single knot component; the bare \code{Diagram}
constructor is not accepted as evidence that those conditions hold.

At crossing $c$, put a tetrahedral cell $C_{c,j}$ in the corner from slot $j$
to slot $j+1$, with subscripts modulo four. Its vertices are the north pole
$0$, the south pole $1$, and the two ends $2,3$ of that corner on the knot.
Only $2,3$ are ideal vertices. Write $\tau=(2\ 3)$, fixing $0,1$, for the
vertex permutation on each face gluing. The complete specification is:
\begin{enumerate}
\item Across slot $j$, glue $C_{c,j-1}$ to $C_{c,j}$ on face
  $1-(j\bmod2)$ by $\tau$. Thus an under-slot joins faces opposite the
  south pole, and an over-slot joins faces opposite the north pole.
\item If an edge pairs slots $(c,j)$ and $(d,k)$, glue face $3$ of
  $C_{c,j}$ to face $2$ of $C_{d,k-1}$ by $\tau$, and face $3$ of
  $C_{d,k}$ to face $2$ of $C_{c,j-1}$ by $\tau$.
\end{enumerate}
All faces occur exactly once. Loops and repeated encounters with a crossing
are allowed: this is a generalized triangulation, not an assertion that all
vertices of its quotient are distinct.

The topological construction is the Weeks crossing-cell decomposition,
documented in the SnapPea kernel.\footnote{Primary source: Jeff Weeks's
description of the triangulation algorithm in
\href{https://github.com/regina-normal/regina/blob/e3a5c7c0e0581706d9e5531b750160b322489950/engine/snappea/kernel/link_complement.cpp}{\texttt{link\_complement.cpp}},
especially the cell decomposition and bigon-collapse sections. The separate
\href{https://github.com/regina-normal/regina/blob/e3a5c7c0e0581706d9e5531b750160b322489950/engine/link/complement.cpp}{Regina implementation}
is an external comparison, not a runtime dependency or copied implementation.}
Here is why its essential hypothesis holds. Thicken the knot near the
equatorial sphere and retain the two polar points. The crossing and
between-crossing meridional cuts partition its exterior into four cells per
crossing. Their two bigon faces collapse along successive under- or over-pass
segments; their other four faces become the displayed tetrahedral faces.
A closed cycle of such bigons would invalidate this collapse. A nonempty
one-component diagram encounters both the under- and over-pass of every
crossing, so no component consists entirely of either kind. The inserted
curl handles the empty case. The polar vertices have sphere links and the
knot vertex has a torus link. The resulting mixed ideal triangulation is
the knot exterior with its torus boundary collapsed to an ideal vertex.
No removal or simplification of the finite poles is required.

The native corner formula uses neither crossing signs nor a chosen strand
orientation. Its correspondence is for the unoriented PL exterior; it does
not supply a marked meridian--longitude basis. The local face prescription,
the bigon-collapse theorem, and the planar source validation are the
mathematical basis of the provenance claim. Agreement with an external
program is additional evidence, not a replacement for that basis.

\subsection{The first compatible finite subdivision}

Give one mixed cell barycentric coordinates
\[
 (x_0,x_1,x_2,x_3),\qquad x_i\ge0,\qquad\sum_i x_i=12.
\]
Retain $x_2\le8$ and $x_3\le8$. These disjoint truncations remove only the
knot vertices. The remaining convex polyhedron has the original poles and
the six directed-edge cut points $E_{vw}$, where $v\in\{2,3\}$, $w\ne v$,
and the two nonzero coordinates of $E_{vw}$ are $x_v=8,x_w=4$.
Its original faces opposite $0,1$ are pentagons and those opposite $2,3$
are quadrilaterals. There are also two triangular truncation faces.

Fan each original face from its interior point $F_f$, whose coordinate $f$
is zero and whose other coordinates are four. Leave each truncation triangle
whole. Cone all boundary triangles to $(3,3,3,3)$. This produces
\[
 5+5+4+4+1+1=20
\]
tetrahedra. Convexity and the face fans prove that their interiors are
disjoint and cover the truncated cell. As an arithmetic control, the sum of
their absolute determinants is
$12^3(1-2/3^3)=1600$; none is degenerate. Volume agreement is a test, not
the only justification of the subdivision.

The old face maps permute barycentric coordinates by $\tau$, so they carry
cut points to cut points and face centres to face centres. The two
subdivisions therefore agree on every paired face. The new truncation
triangles are precisely the finite boundary. Keeping the finite poles fills
their spherical links; no extra spherical boundary components are created.

\subsection{Pulling from a retained pole: four times fewer tetrahedra}

The first implementation unnecessarily inserts a centre in every old face
and every crossing cell. Those centres make an arbitrary face permutation
easy to accommodate, but every actual gluing fixes both poles. This extra
structure permits a smaller compatible subdivision.

Write $P_0,P_1$ for the poles. Pull the polyhedron from $P_0$: triangulate
the facets not containing $P_0$ and cone them to $P_0$. The facets in question
are the pentagon opposite $0$ and the two truncation triangles. Fan the
pentagon from $P_1$. The complete tetrahedron list becomes
\[
\begin{split}
 &(P_0,P_1,E_{21},E_{23}),\quad (P_0,P_1,E_{23},E_{32}),\\
 &(P_0,P_1,E_{32},E_{31}),\quad (P_0,E_{20},E_{21},E_{23}),\\
 &(P_0,E_{30},E_{31},E_{32}).
\end{split}
\]
Convexity proves coverage: a ray from $P_0$ through an interior point exits
through exactly one facet not containing $P_0$. Distinct facet triangles
give disjoint tetrahedral interiors. On each old face, the induced
triangulation is the fan from its smallest numbered finite pole. The map
$\tau$ fixes both poles, so these fans agree across every gluing, including
the identifications between old faces $2$ and $3$. The two cut faces remain
single triangles. There are now five nondegenerate tetrahedra, twelve
external triangles and four internal face pairs per crossing cell. Their
absolute determinant sum is again $1600$.

\begin{proposition}
For every validated classical one-component PD, the maintained constructor
returns a compact triangulation of its unoriented knot exterior with exactly
$20N$ tetrahedra and $8N$ boundary triangles. The original centred variant
uses $80N$ tetrahedra with the same boundary count. The only boundary component
is a torus. Its output format is accepted by the native normal-surface geometry
validator.
\end{proposition}
\begin{proof}
Apply the crossing-cell theorem and the pulling subdivision to its $4N$
cells. Truncating the unique ideal knot vertex recovers a regular compact
exterior. Each cell contributes five tetrahedra and two boundary triangles.
The previous centred construction instead contributes twenty tetrahedra
per cell. All other faces are paired by the local subdivisions or original
face identifications.
\end{proof}

The default API now selects \code{pulling}; \code{subdivision='centred'}
retains the original output byte for byte. The independent checker accepts
both, distinguishing them by the tetrahedron count relative to the validated
source. No arbitrary retriangulation is accepted. Standard normal-coordinate
vectors for the default output have $140N$ entries rather than $560N$.
This is a factor-four reduction in explicit geometry and coordinate dimension,
not an asymptotic improvement to a complete recognition algorithm.

\subsection{Independent replay and its scope}

The producer labels points symbolically as poles and directed-edge cut points,
with centres only in the legacy variant. It pairs equal local facets and
transports symbolic labels across old faces. The checker does not import that module.
It uses a separately specified list of integer-coordinate polygons, derives
the old face from the coordinate that vanishes on all three facet points,
and recovers the neighbouring corner directly from PD parity and edge mates.
It checks every target tetrahedron and every vertex permutation. Its strict
schema rejects booleans masquerading as indices, partial pairings and extra
fields. Cancellation callbacks propagate, including false-valued callables.

The checker accepts canonical numbering only. An isomorphic retriangulation
requires additional evidence; even a valid one-torus-boundary manifold cannot
be substituted for the source exterior. Both paths revalidate the source PD.
Once this geometry check succeeds, an independently accepted compressing-disc
certificate in the very same triangulation proves that the source knot is
unknotted: a knot exterior in $S^3$ is irreducible, and a compression of its
torus boundary makes it a solid torus. The construction does not itself
produce such a disc. In particular, a boundary vertex-linking disc has
inessential boundary and supplies no positive unknot certificate.

After the normalized PD and its edge mates are available, construction and
replay perform $O(N)$ fixed-size local combinatorial operations. The explicit
output occupies $O(N\log(N+1))$ bits because its face records contain global
tetrahedron indices. Source validation is polynomial preprocessing; this is
not a claim of linear bit complexity for arbitrary input labels or of cheap
normal-surface search on the resulting triangulation. The explicit geometry
is still substantially larger than typical simplified external-engine
triangulations. The API is optional and no default recognition path changes.

\subsection{Source controls, compatibility and complete-operation costs}

The final native audit uses the nineteen-case maintained corpus, the empty
circle and the valid one-component closures retained from 240 seeded random
braids. These give 84 source diagrams. Each is tested in its original form,
as a mirror, and after reversing the crossing list and rotating each PD row
by two slots: 252 diagram variants, with at most 141 crossings.

Every mixed crossing triangulation is isomorphic to Regina's unsimplified
diagram complement and has vertex-link Euler characteristics $2,2,0$.
For each variant, both finite subdivisions pass the native manifold validator
and Regina's finite orientable one-torus-boundary checks: 504 comparisons.
The current source checker accepts both outputs, and the original checker
accepts the retained centred output: 756 successful source replays. All 252
centred outputs exactly equal the first published constructor at commit
\code{8354e8aab}. The external engine is Regina 7.4, distribution 7.4.1.
Neither engine's manifold checks alone would establish knot provenance.

The focused tests disable the producer during independent replay, check
nondegeneracy and exact volume for both local subdivisions, reject all 400
single-face mutations of the two one-crossing outputs, and reject a
reciprocal but incorrect gluing. They also reject a trefoil exterior offered
for a three-crossing unknot source, malformed schemas and unchecked invalid
PDs. Cancellation propagates without input mutation. A source-bound boundary
vertex-link disc correctly has zero essential-disc count; changing that
certificate's count to one is rejected. Seven small external solid-torus
controls cover the empty circle, curls of both signs, a reducible
three-crossing unknot, a trefoil, a figure-eight and a stabilized circle.
The original implementation passed 1,149 maintained tests. The final version
adds the retained-subdivision contract test and passes all 1,150 maintained
tests in 216.795 seconds. Its eight focused tests pass in 1.353 seconds.
The final geometry audit took 56.440 seconds while the suite also ran;
that elapsed time is not used as a performance comparison.

There was no newly discovered normal meridian at this construction-only
checkpoint; Section~\ref{sec:cocycle-seeds} subsequently supplies native
candidates and source-bound disc proofs. Two
exploratory external calls to \code{nonTrivialSphereOrDisc()}, one on the
80-tetrahedron circle and one on the 20-tetrahedron circle, each reached its
20-second process cap without returning a surface. These are incomplete
searches, not negative results or a speed comparison. The existing supplied
normal-vector algorithms can now be applied to source-authenticated geometry,
but obtaining useful vectors efficiently remains an independent obligation.

\subsubsection{Measured effect of the pulling subdivision}

The final benchmark runs serially after the suite, audit and publication of
the runtime change. It compares the pinned first constructor and checker at
\code{8354e8aab} against the final pulling implementation, using the same
maintained manifold validator. Each timed call starts afresh, constructs the
exterior, independently replays its source geometry and validates the finite
manifold. Serialization for the stored output hash is outside the timer.
Both variants perform the complete stated operation; they have different
canonical subdivisions and therefore different expected output hashes.

The four inputs are stabilized circles with $1,8,32,128$ crossings. There
are five rounds and four interleaved arms: old, old repeat, new and new
repeat, plus one warmup per arm and input. All eighty measured calls and
sixteen warmups complete in a total driver elapsed time of 24.675 seconds.
Hashes agree within each subdivision across every repeat. The table gives
unpaired time medians and medians of paired ratios.

\begin{center}
\small
\begin{tabular}{rrrrrrr}
\toprule
Crossings & New tets & Old (ms) & New (ms) & Old/new & Old A/A & New A/A \\
\midrule
1 & 20 & 8.444 & 2.269 & 3.748 & 0.996 & 1.043 \\
8 & 160 & 73.140 & 17.259 & 4.346 & 1.032 & 0.997 \\
32 & 640 & 285.732 & 67.210 & 4.228 & 1.021 & 0.998 \\
128 & 2560 & 1217.252 & 284.316 & 4.299 & 0.998 & 0.996 \\
\bottomrule
\end{tabular}

\end{center}

The complete-call old/new ratios are $3.748$--$4.346$. Old A/A ratios are
$0.996$--$1.032$ and new A/A ratios are $0.996$--$1.043$. At 128 crossings,
the median falls from $1.217252$ seconds to $0.284316$ seconds; tetrahedra
fall from 10,240 to 2,560. This is consistent with the exact factor-four
reduction in geometry. It measures construction, source replay and manifold
validation, not a normal-vector search, a simplified minimal exterior or
whole recognition. It supplies no new general quasipolynomial theorem.


Reproduction uses \path{../fast/normal_orbit_research/exteriors.py}, with
\code{audit} and \code{benchmark} modes. The driver pins all 457 Python and
fixture sources before and after each run, and separately pins the two
baseline modules loaded from Git. Final records have the
\code{diagram-exterior-} prefix in \path{data/}; the first implementation's
audit, suite and A/A timing records retain the
\code{diagram-exterior-centred-} prefix. The coordinate and gluing proofs are
in the maintained article, not delegated to the optional external controls.

